\documentclass[11pt]{article}
\usepackage[margin=1in]{geometry}
\usepackage{amsmath,amssymb,amsthm}
\usepackage{graphicx}
\usepackage[colorlinks=true,linkcolor=blue,urlcolor=blue,citecolor=blue]{hyperref}
\setlength{\parindent}{0pt}
\setlength{\parskip}{0.65em}

\newtheorem{theorem}{Theorem}
\newtheorem{lemma}{Lemma}
\newcommand{\A}{\mathcal A}
\newcommand{\B}{\mathcal B}
\newcommand{\Hh}{\mathcal H}
\newcommand{\K}{\mathcal K}
\newcommand{\BH}{\mathcal B(\Hh)}
\newcommand{\UCP}{\operatorname{UCP}}

\title{The nonseparable $C^*$-Krein--Milman theorem\\
\large Finite BW localization and pure extensions}
\author{Run \texttt{fa\_banach\_001} --- \texttt{agent\_lane\_11}}
\date{13 August 2026}

\begin{document}
\maketitle

\textbf{Status.} \texttt{candidate\_full\_solution\_likely\_valid}.
The argument answers one explicit open question in \cite{BK}; it should receive
expert review before being treated as established literature.

\section*{The source question}

For a unital $C^*$-algebra $\A$ and a separable Hilbert space $\Hh$, write
\[
 S_{\Hh}(\A)=\UCP(\A,\BH).
\]
The bounded weak (BW) topology is pointwise weak-operator convergence. Bhat and
Kumar prove that $S_{\Hh}(\A)$ is the BW-closure of the $C^*$-convex hull of its
$C^*$-extreme points when $\A$ is separable or a type I factor. On source-PDF
page 19 they state that the result is unknown for general nonseparable
$C^*$-algebras \cite[Section 5]{BK}.

\begin{center}
\includegraphics[width=0.97\textwidth]{figures/open_problem_crop.png}
\end{center}

\section*{Proof intuition}

BW approximation is local: a basic neighborhood tests only finitely many elements
of $\A$, and those elements generate a separable subalgebra $\B$. The source theorem
therefore supplies a $C^*$-convex approximation on $\B$.

The subtle point is extension. An arbitrary $C^*$-extreme map on $\B$ need not be
known to have a $C^*$-extreme extension. Fortunately, the source proof uses much more
structured extreme maps. Each is built from a pure UCP map and a one-dimensional
pure compression of that same map. Pure UCP maps have pure extensions. Extending the
first map and taking the same vector compression recreates the structured block on
$\A$, where the source's direct-sum theorem again makes it $C^*$-extreme. Thus the
same operator coefficients lift the local approximation to all of $\A$.

\section*{Two extension lemmas}

Recall that a completely positive map $\alpha$ is \emph{pure} if
$0\leq\gamma\leq\alpha$ implies $\gamma=t\alpha$ for some $t\in[0,1]$.

\begin{lemma}[pure extension]\label{lem:pure-extension}
Let $\B\subseteq\A$ be unital $C^*$-algebras and let
$\alpha:\B\to\mathcal B(\K)$ be pure and unital completely positive. Then
$\alpha$ has a pure UCP extension
$\widetilde\alpha:\A\to\mathcal B(\K)$.
\end{lemma}

\begin{proof}
By Arveson's extension theorem \cite{Arveson}, the set
\[
 \mathcal F=\{\beta\in\UCP(\A,\mathcal B(\K)):\beta|_{\B}=\alpha\}
\]
is nonempty. It is BW-compact and convex. It is also a face of
$\UCP(\A,\mathcal B(\K))$: if
$t\beta_1+(1-t)\beta_2\in\mathcal F$, with $0<t<1$, then
$t\beta_1|_{\B}\leq\alpha$. Purity gives
$t\beta_1|_{\B}=s\alpha$; evaluation at $1$ gives $s=t$, hence
$\beta_1|_{\B}=\alpha$, and similarly for $\beta_2$.

Choose an extreme point $\widetilde\alpha$ of $\mathcal F$. Since $\mathcal F$
is a face, it is linearly extreme in the full UCP set. Suppose
$0\leq\gamma\leq\widetilde\alpha$. Purity of the restriction gives
$\gamma|_{\B}=s\alpha$. Thus $\gamma(1)=sI$. If $0<s<1$, then
\[
 \widetilde\alpha
 =s\,(s^{-1}\gamma)+(1-s)\bigl((1-s)^{-1}(\widetilde\alpha-\gamma)\bigr)
\]
is a convex decomposition into UCP maps. Linear extremality forces
$\gamma=s\widetilde\alpha$. The cases $s=0,1$ follow because a CP map vanishing
at the unit is zero. Hence $\widetilde\alpha$ is pure.
\end{proof}

The next lemma isolates precisely the blocks appearing in the proof of
\cite[Lemma 5.1]{BK}.

\begin{lemma}[lifting the source blocks]\label{lem:block}
Let $\K\subseteq\Hh$, let $\alpha:\B\to\mathcal B(\K)$ be pure UCP, and let
$h\in\K$ be a unit vector. Put
\[
 \vartheta(b)=\langle h,\alpha(b)h\rangle,
 \qquad
 \xi(b)=\alpha(b)\oplus\vartheta(b)I_{\K^\perp}.
\]
Assume (as for the source blocks) that $\vartheta$ is a pure state. Then $\xi$
has a $C^*$-extreme UCP extension
$\widetilde\xi:\A\to\BH$.
\end{lemma}

\begin{proof}
Take a pure extension $\widetilde\alpha$ from Lemma~\ref{lem:pure-extension} and
define
$\widetilde\vartheta(a)=\langle h,\widetilde\alpha(a)h\rangle$.
The minimal Stinespring representation of a pure CP map is irreducible
\cite{Arveson}. Every nonzero vector in an irreducible representation is cyclic;
therefore $\widetilde\vartheta$ is a pure state and extends $\vartheta$.

Set
\[
 \widetilde\xi(a)=\widetilde\alpha(a)
       \oplus\widetilde\vartheta(a)I_{\K^\perp}.
\]
The scalar summand is a compression of the same irreducible Stinespring
representation as $\widetilde\alpha$. Their Stinespring ranges are nested, and the
one-type multiplicity nest has countable completion. The direct-sum criterion
\cite[Theorem 3.6]{BK} therefore makes $\widetilde\xi$ $C^*$-extreme. Plainly
$\widetilde\xi|_{\B}=\xi$.
\end{proof}

The same lemma includes the fixed residual block needed below: choose a pure state
$\omega$ of $\B$. Then $b\mapsto\omega(b)I_{\Hh}$ is a one-type amplification,
hence $C^*$-extreme by the same direct-sum criterion, and a pure extension of
$\omega$ gives a $C^*$-extreme extension of this block.

\section*{Main theorem}

\begin{theorem}\label{thm:main}
Let $\A$ be any unital $C^*$-algebra and let $\Hh$ be separable. Then
\[
 \boxed{
 S_{\Hh}(\A)=
 \overline{\operatorname{co}_{C^*}
 \bigl(\operatorname{Cext}S_{\Hh}(\A)\bigr)}^{\mathrm{BW}}.}
\]
Thus the full-generality question in \cite{BK} has an affirmative answer.
\end{theorem}

\begin{proof}
Fix $\phi\in S_{\Hh}(\A)$ and a basic BW neighborhood $\mathcal N$ of $\phi$.
It is enough to consider finitely many tests
\[
 |\langle(\psi(a_r)-\phi(a_r))u_s,v_s\rangle|<\varepsilon
 \quad(1\leq r\leq m,\ 1\leq s\leq q).
\]
Let $\B=C^*(1,a_1,\ldots,a_m)$. This is a separable unital $C^*$-subalgebra,
so the separable case of \cite[Theorem 5.3]{BK} applies to
$\phi_0=\phi|_{\B}$.

We use one feature visible in the proof of that theorem. Its Case I first
approximates $\phi_0$ by a compression of a countable direct sum of irreducible
representations. That compression is a WOT-sum of pure CP maps. The proof of
\cite[Lemma 5.1]{BK} then approximates it by finite $C^*$-convex combinations
\[
 \eta=\sum_{j=0}^{N}T_j^*\xi_j(\mathord\cdot)T_j,
 \qquad \sum_{j=0}^{N}T_j^*T_j=I_{\Hh},                         \tag{1}
\]
where every non-residual $\xi_j$ is, up to a unitary, exactly a block from
Lemma~\ref{lem:block}. The residual coefficient tends strongly to zero and its
map may be chosen as $\omega(\mathord\cdot)I_{\Hh}$ for a pure state $\omega$
of $\B$. Choosing the two approximation stages sufficiently accurately gives
$\eta$ satisfying all the tests defining $\mathcal N$.

By Lemma~\ref{lem:block}, extend each $\xi_j$ to a $C^*$-extreme
$\widetilde\xi_j\in S_{\Hh}(\A)$; use the pure-state amplification extension for
the residual block. Keep the coefficients in (1) and define
\[
 \widetilde\eta=\sum_{j=0}^{N}T_j^*\widetilde\xi_j(\mathord\cdot)T_j.
\]
This belongs to the $C^*$-convex hull of the $C^*$-extreme points of
$S_{\Hh}(\A)$, and $\widetilde\eta|_{\B}=\eta$. Since every tested $a_r$ lies in
$\B$, the map $\widetilde\eta$ belongs to $\mathcal N$.

Every basic BW neighborhood of every $\phi$ therefore meets the stated hull,
which is precisely the asserted BW-closure identity.
\end{proof}

\section*{Verification, limitations, and novelty check}

The proof deliberately does not assert that arbitrary $C^*$-extreme maps extend.
It lifts only the pure/compression blocks actually constructed in the source's
Lemma 5.1. The compact extension-fiber argument uses nets and needs no
metrizability. The conclusion is local BW closure, so no diagonal sequence over
all of a nonseparable algebra is required.

The result retains the source convention that $\Hh$ is separable. It removes the
separability and type I assumptions on $\A$; it does not address the paper's other
questions about the structure of normal $C^*$-extreme maps or whether every
$C^*$-extreme map is linearly extreme.

A bounded novelty search on 13 August 2026 covered the run indexes, the parsed
arXiv corpus, exact-title citations, and primary-source searches for the full
nonseparable BW theorem and the pure-extension route. It found later work on
specialized classes but no statement of Theorem~\ref{thm:main}. This supports
candidate novelty but is not an exhaustive priority claim.

\begin{thebibliography}{9}
\bibitem{Arveson}
W. Arveson, \emph{Subalgebras of $C^*$-algebras}, Acta Math. 123 (1969),
141--224. In particular, the CP extension and pure-extension machinery.

\bibitem{BK}
B. V. Rajarama Bhat and M. Kumar, \emph{$C^*$-extreme maps and nests},
arXiv:2103.09600; J. Funct. Anal. 282 (2022), Paper No. 109397.
\end{thebibliography}

\end{document}
